% Источники: exam/materials/преподаватель/2020/27-04-2020-Model_L10_3_.pdf; exam/materials/моделирование.pdf, раздел 34.
\section{Коэффициент корреляции Пирсона.}

\subsection*{Коэффициент корреляции Пирсона}

Применяется при наличии \textbf{линейной стохастической зависимости} между случайными величинами $X$ и $Y$.

\[
  r_{XY} = \frac{\sigma_{XY}}{\sigma_X \cdot \sigma_Y},
\]
где
\[
  \sigma_{XY} = \frac{1}{n}\sum_{i=1}^n (X_i - \bar{X})(Y_i - \bar{Y})
\]
— выборочная ковариация, а
\[
  \sigma_X = \sqrt{\frac{1}{n}\sum_{i=1}^n (X_i - \bar{X})^2}, \quad
  \sigma_Y = \sqrt{\frac{1}{n}\sum_{i=1}^n (Y_i - \bar{Y})^2}
\]
— выборочные среднеквадратические отклонения.

\subsection*{Алгоритм}

\begin{enumerate}
  \item Вычислить выборочные средние $\bar{X}$ и $\bar{Y}$.
  \item Вычислить ковариацию $\sigma_{XY}$.
  \item Вычислить $\sigma_X$ и $\sigma_Y$.
  \item Найти $r_{XY} = \sigma_{XY}/(\sigma_X\sigma_Y)$.
\end{enumerate}

\subsection*{Интерпретация}

$r_{XY} \in [-1, 1]$: близкое к $\pm1$ — сильная линейная зависимость, близкое к $0$ — зависимость отсутствует или нелинейна.

\subsection*{Пример}

\begin{center}
\begin{tabular}{ccccc}
  \hline
  $i$ & $X_i$ & $Y_i$ & $(X_i-\bar{X})$ & $(Y_i-\bar{Y})$ \\
  \hline
  1 & 1 & 2 & $-2$ & $-2$ \\
  2 & 2 & 3 & $-1$ & $-1$ \\
  3 & 3 & 4 & $0$ & $0$ \\
  4 & 4 & 5 & $+1$ & $+1$ \\
  5 & 5 & 6 & $+2$ & $+2$ \\
  \hline
\end{tabular}
\end{center}

$\bar{X}=3$, $\bar{Y}=4$;
$\sigma_{XY} = \frac{1}{5}(4+1+0+1+4)=2$;
$\sigma_X=\sigma_Y=\sqrt{2}$.

\[
  r_{XY} = \frac{2}{\sqrt{2}\cdot\sqrt{2}} = \frac{2}{2} = 1.
\]

Вывод: полная прямая линейная зависимость.

\clearpage
